
\chapter*{ЗАКЛЮЧЕНИЕ}
\addcontentsline{toc}{chapter}{ЗАКЛЮЧЕНИЕ}

В результате выполнения данной дипломной работы была достигнута цель, а именно разработан метод обеспечения дифференциальной приватности при генерации синтетических данных. 

Для достижения поставленной цели были решены  следующие задачи:
\begin{itemize}[label=---]
	\item проанализировать и классифицировать существующие методы генерации синтетических данных и дифференциальной приватности;
	\item сформулировать концепцию комбинированного метода на уровне идеи, формализовать задачу в виде функциональной модели;
	\item выделить ключевые этапы разрабатываемого метода и описать их в виде детализированной диаграммы IDEF0 и схем алгоритмов;
	\item разработать программное обеспечение, реализующее разработанный метод;
	\item исследовать характеристики разработанного метода при различных входных
	данных;
	\item провести сравнение результатов работы реализованного метода с альтернативными методами.
	
\end{itemize}
